\documentclass[11pt]{article}
\usepackage[margin=1in]{geometry}
\usepackage{amsmath,amssymb,amsthm,mathtools}
\usepackage[colorlinks=true,urlcolor=blue,linkcolor=blue]{hyperref}
\newtheorem{proposition}{Proposition}
\newcommand{\tr}{\operatorname{tr}}
\newcommand{\str}{\operatorname{str}}
\newcommand{\LL}{\mathcal L}
\newcommand{\Res}{\operatorname{Res}}
\title{Ordinary specialization and its geometric nilpotent residue}
\author{Focused independent audit}
\date{11 October 2026}
\begin{document}
\maketitle

\paragraph{Finding.}
Under the divisor hypotheses in
\emph{holomorphic\_continuity\_crossing.tex}, specializing the spectral
fiber at a crossing gives exactly $Z(0)=G(0)c e_1$. Neither ordinary
parameter continuity nor joint holomorphy follows from those hypotheses.
The additional calculation here identifies the possible discontinuity
with a nilpotent trace in the published contact variation formula.
Unitary Stokes does not annihilate that trace: it reduces it to the
same weighted $du$--$du_*$ pairing that governs conformal first variation.
This note does not independently verify the geometric existence of the
claimed crossing.

\section{The two values and the missing cancellation}
Work locally at a crossing and write a deformation parameter as $u$,
with crossing at $u=0$. Suppose
\[
 Z(s,u)=G(s,u)\frac{(s-c(u))(s^2-e_1(u)s+e_2(u))}{s-t(u)},
 \qquad c(0)\ne0,\quad t(0)=e_2(0)=0.
\]
The fixed-cut remainder $G$ is holomorphic and nonzero in $s$.
Specializing $u=0$ gives, as a meromorphic function of $s$,
\[
 Z(s,0)=G(s,0)(s-c(0))(s-e_1(0)).
\]
Thus its ordinary value is $G(0,0)c(0)e_1(0)$, including zero if
$e_1(0)=0$.

For comparison assume a transverse, sufficiently regular local
deformation, so $v:=t'(0)\ne0$, and put $e_2=tH$ with
$H(0)=e_2'(0)/v$. If ordinary nonresonant values have constant value
$A\ne0$ on one side and $G,c$ are continuous, then
\[
 H(0)=\frac{A}{G(0,0)c(0)}\ne0.
\]
The continued nonresonant value is $G(0,0)c(0)H(0)$.
It equals the ordinary crossing value precisely when $H(0)=e_1(0)$.
These are different operations on a meromorphic family.

For holomorphic parameter germs, monic division gives
\begin{equation}\label{division}
 \frac{s^2-e_1s+tH}{s-t}
 =s+t-e_1+\frac{t(t-e_1+H)}{s-t}.
\end{equation}
Since $G(s,u)(s-c(u))$ is a unit near the crossing, joint holomorphy
requires the germ identity $H=e_1-t$. Fiberwise cancellation at $u=0$
does not imply this identity. If $H(0)\ne e_1(0)$, paths $s=a t(u)$,
$a\ne1$, have the limiting values
\[
 \lim_{u\to0}Z(a t(u),u)
 =G(0,0)c(0)\frac{a e_1(0)-H(0)}{a-1}.
\]
Even equality at the single point is weaker than joint holomorphy:
unless the remainder in \eqref{division} vanishes as a germ, nearby
genuine poles can be approached along curves close to $s=t(u)$.

The primary references distinguish the two objects. DGRS Theorem~2
only establishes local constancy on the set with no zero resonance.
Chaubet--Dang (6.4)--(6.5) removes the cut divisor, while (6.7) retains
the zero-complex refined torsion in the relation to the specialized
Laurent coefficient \cite{DGRS,CD}. Neither statement identifies
$H(0)$ with $e_1(0)$.

\section{An exact Laurent diagnostic from contact variation}
Let $L=\LL_X^{d_\theta}$, so the spectral matrix is $B=-L$, and put
$K=\Pi\alpha\wedge\iota_{\dot X}$ on the finite full cut.
In flow dimension five, CD Lemma~7.4 gives for its horizontal
alternating divisor $F=Z/G$
\begin{equation}\label{variation}
 \partial_u\log F(s,u)
 =-\str_{C_{\mathrm{cut}}^\bullet}
       \bigl(K L(L+s)^{-1}\bigr).
\end{equation}
The nonzero blocks contribute holomorphic terms near $s=0$.
At the crossing, direct differentiation of the displayed divisor gives
\[
 \partial_u\log F(s,0)
 =-\frac{c'}{s-c}
   +\frac{-e_1's+e_2'}{s(s-e_1)}+\frac{v}{s}.
\]
If $m=2$, then $e_1(0)\ne0$ and $L^2=0$ on the full zero block.
Comparison of the $s^{-1}$ coefficients in \eqref{variation} yields
\begin{equation}\label{m2trace}
 \str_{C^\bullet(0)}(KL)
    =v\left(\frac{H(0)}{e_1(0)}-1\right).
\end{equation}
Thus a transverse crossing has the continued ordinary value precisely
when this trace vanishes. In particular, nilpotence by itself does not
give the required conclusion.

If $m=3$, then $e_1(0)=0$ and $L^3=0$ on the zero block. Since
$L(L+s)^{-1}=L/s-L^2/s^2$, comparison gives
\begin{equation}\label{m3trace}
 \str_{C^\bullet(0)}(KL^2)=e_2'(0)=vH(0),\qquad
 \str_{C^\bullet(0)}(KL)=e_1'(0)-v.
\end{equation}
For a transverse crossing with nonzero continued value, the first
trace is nonzero. The $m=3$ fiber itself has a simple zero.

\section{The geometric Stokes reduction of the traces}
We now work on the actual outgoing and incoming zero spaces, using
the Hermitian flat-bundle pairing. At a crossing their horizontal
dimensions are $(0,1,m,1,0)$. Choose nonzero outgoing and incoming
horizontal zero one-forms $u,u_*$, and write
\[
 x=d_\theta u=Du,\qquad x_*=d_\theta u_*=Du_*,\qquad
 W=d\alpha\wedge.
\]
Full exactness makes $x,x_*$ nonzero and $D:V^2\to V^3$ surjective.
The outgoing--incoming pairing and Lefschetz give
\[
 B_1:=\int_M\alpha\wedge d\alpha\wedge
                  \langle u\wedge u_*\rangle_E\ne0.
\]
Choose $y\in V^2$ with $Dy=Wu$. For two horizontal degree-two
states use
\[
 P(a,a_*)=\int_M\alpha\wedge\langle a\wedge a_*\rangle_E.
\]
Unitarity permits Stokes with the twisted differential. The horizontal
five-form term vanishes pointwise, giving
\begin{equation}\label{stokes}
 P(y,x_*)=-B_1,\qquad P(x,x_*)=0,\qquad
 P(z,x_*)=0\quad(z\in\ker D).
\end{equation}
Consequently $\ell_y(a):=-P(a,x_*)/B_1$ detects the $y$ coordinate.
Pairing with an incoming zero state also removes the outgoing
projector, so this formula applies to compressed smooth multiplication.

Set $h=\alpha(\dot X)$, and define
\[
 S_h:=P(hx,x_*)
      =\int_M h\alpha\wedge\langle d_\theta u\wedge
                                             d_\theta u_*\rangle_E.
\]
Contact splitting shows for every $j\ge1$
\begin{equation}\label{splittrace}
 \str_{C^\bullet(0)}(K L^j)
       =-\tr_{V^2}\bigl((\Pi_{V^2}h)L^j\bigr),
\end{equation}
because $L=0$ on $V^1,V^3$ and the diagonal action of $K$ on
$\alpha\wedge V^2$ is multiplication by $h$.

For $m=2$, $\ker D=\mathbb Cx$, and $[B,D]=0$ gives
$By=\beta x$, $Bx=0$. Hence
\[
 \tr_{V^2}(hB)=\beta\ell_y(hx)=-\beta S_h/B_1,
 \qquad \str_{C^\bullet(0)}(KL)=-\beta S_h/B_1.
\]
For $m=3$, choose $z\in\ker D$ complementary to $x$. Nilpotence
and $[B,D]=0$ give constants $\delta,\eta,\epsilon$ with
\[
 Bx=0,\quad Bz=\delta x,\quad By=\epsilon x+\eta z,
 \quad B^2y=\delta\eta x.
\]
Thus
\begin{equation}\label{rankthree}
 \str_{C^\bullet(0)}(KL^2)=\delta\eta S_h/B_1.
\end{equation}
Full exactness in this case actually forces $\delta\eta\ne0$:
in degree two the images of $z,y,\alpha\wedge u$ are
$-\delta\alpha\wedge x$,
$Wu-\alpha\wedge(\epsilon x+\eta z)$, and
$Wu-\alpha\wedge x$. They must have rank three, which is
equivalent to $\delta\eta\ne0$. The nilpotent block therefore has
length three.

\section{Conformal first variation does not kill the weighted term}
For the conformal family $g_u=e^{-2u b}g$ under radial identification,
$\dot\alpha=-b\alpha$, so differentiating $\alpha(X)=1$ gives
$h=\alpha(\dot X)=b\circ\pi$. CDDP Lemma~4.2, (4.22), and (4.38)
give the one-form carrier variation
\begin{equation}\label{carrier}
                 v=t'(0)=S_h/B_1.
\end{equation}
At a simple unitary twisted one-zero block the same proof works:
the connection preserves the coefficient pairing, Cartan holds, and
$W d_\theta u=W d_\theta u_*=0$ because the corresponding horizontal
degree-four zero space vanishes. These are precisely the identities
needed in the integration-by-parts calculation.

Combining \eqref{carrier} with the preceding geometric formulas gives
\[
 m=2:\quad \str(KL)=-\beta v,
 \qquad
 m=3:\quad \str(KL^2)=\delta\eta v.
\]
At a transverse crossing, \eqref{m2trace} identifies
$\beta=1-H(0)/e_1(0)$, and \eqref{m3trace} identifies
$\delta\eta=H(0)$. In the $m=2$ case ordinary continuity is therefore
equivalent to semisimplicity of the middle zero block. In the $m=3$
case full exactness and transversality forbid the vanishing proposed
for the quadratic nilpotent trace.

The concrete incorrect Stokes inference would be to treat $h x$ as
exact and conclude $S_h=0$. Direct Stokes instead gives
\begin{equation}\label{weightedstokes}
 S_h=\int_M dh\wedge\alpha\wedge
                     \langle u\wedge d_\theta u_*\rangle_E.
\end{equation}
The term with $d\alpha$ is a horizontal five-form and vanishes, but
the $dh$ term remains. Multiplication satisfies
$h d_\theta u=d_\theta(hu)-dh\wedge u$.
For nonconstant base $b$ this is not an exactness-preserving operation.
CDDP's conformal nondegeneracy calculation is built from this very
weighted pairing and supplies nonzero geometric examples at the
hyperbolic endpoint. It does not assert that it vanishes at acyclic
crossings \cite{CDDP}.

Likewise the finite trace of $[L,\Pi h\Pi]$ vanishes by cyclicity,
but this only concerns the compressed commutator. It is not an
identity setting $\tr(\Pi h\Pi L)$ to zero. Incoming Stokes duality
and geodesic reversal give an indefinite pairing; they do not turn
this nilpotent trace into a positive-definite skew-adjoint trace.

\paragraph{Audit conclusion.}
The ordinary spectral cancellation is valid. A rescue by continuous
nonresonant values would require the extra equality $H(0)=e_1(0)$;
joint holomorphy would require the stronger germ identity $H=e_1-t$.
At a transverse $m=2$ crossing, the needed equality is exactly
middle-degree semisimplicity. The actual contact and unitary Stokes
calculation reduces its obstruction to a weighted conformal pairing,
rather than eliminating it. No invalid inference in the supplied
fiber-specialization step has been found.

\begin{thebibliography}{9}
\bibitem{CD} Y. Chaubet and N. V. Dang,
\emph{Dynamical torsion for contact Anosov flows},
\url{https://arxiv.org/html/1911.09931v2},
(6.4)--(6.7), Lemma 7.4 and (7.7)--(7.10).
\bibitem{DGRS} N. V. Dang, C. Guillarmou, G. Rivi\`ere and S. Shen,
\emph{The Fried conjecture in small dimensions},
\url{https://webusers.imj-prg.fr/~shu.shen/publication/Invention_DGRS.pdf},
Theorem 2 and Proposition 7.3.
\bibitem{CDDP} M. Ceki\'c, B. Delarue, S. Dyatlov and G. P. Paternain,
\emph{The Ruelle zeta function at zero for nearly hyperbolic
3-manifolds}, \url{https://arxiv.org/html/2009.08558},
Lemma 4.2, (4.22), (4.38), and Theorem 5.
\end{thebibliography}
\end{document}
